
\documentclass{article}
\usepackage{colortbl}
\usepackage{makecell}
\usepackage{multirow}
\usepackage{supertabular}

\begin{document}

\newcounter{utterance}

\centering \large Interaction Transcript for game `matchit\_ascii', experiment `same\_grid', episode 1 with nemotron{-}3{-}ultra.
\vspace{24pt}

{ \footnotesize  \setcounter{utterance}{1}
\setlength{\tabcolsep}{0pt}
\begin{supertabular}{c@{$\;$}|p{.15\linewidth}@{}p{.15\linewidth}p{.15\linewidth}p{.15\linewidth}p{.15\linewidth}p{.15\linewidth}}
    \# & $\;$A & \multicolumn{4}{c}{Game Master} & $\;\:$B\\
    \hline

    \theutterance \stepcounter{utterance}  
    & & \multicolumn{4}{p{0.6\linewidth}}{
        \cellcolor[rgb]{0.9,0.9,0.9}{
            \makecell[{{p{\linewidth}}}]{
                \texttt{\tiny{[P1$\langle$GM]}}
                \texttt{Tá tú ag glacadh páirte i gcluiche tomhais comhoibríoch.} \\
\texttt{Is é an sprioc a fháil amach an bhfuil an greille seo agus pictiúr eile nach féidir liomsa amháin a fheiceáil, mar an gcéanna.} \\
\texttt{Déan cur síos ar do ghreille ar dtús, le do thoil. Ansin, tabharfaidh mé mo chur síos agus is féidir linn ceisteanna a chur ar a chéile faoi na greillí chun a dhéanamh amach an bhfuil siad mar an gcéanna. Anois cuir tús le do chur síos gairid ar an ngreille le "TUAIRISC:" agus an cur síos ina dhiaidh. Ná cuir aon rud eile leis. Seo é do ghreille:} \\
\texttt{▢ ▢ X ▢ ▢} \\
\texttt{▢ ▢ X ▢ ▢} \\
\texttt{▢ ▢ X ▢ ▢} \\
\texttt{▢ ▢ X ▢ ▢} \\
\texttt{X X X X X} \\
            }
        }
    }
    & \\ \\

    \theutterance \stepcounter{utterance}  
    & \multicolumn{4}{p{0.6\linewidth}}{
        \cellcolor[rgb]{0.9,0.9,0.9}{
            \makecell[{{p{\linewidth}}}]{
                \texttt{\tiny{[P1$\rangle$GM]}}
                \texttt{TUAIRISC: Greille 5x5. Trídchló ghearrthóir (colún 3) agus líne bhunúsach (ráille 5) atá líonta le X, ag cruthú cruth T. Na cealla eile folamh.} \\
            }
        }
    }
    & & \\ \\

    \theutterance \stepcounter{utterance}  
    & & & \multicolumn{2}{p{0.3\linewidth}}{
        \cellcolor[rgb]{0.9,0.9,0.9}{
            \makecell[{{p{\linewidth}}}]{
                \texttt{\tiny{[GM$|$GM]}}
                \texttt{continue} \\
            }
        }
    }
    & & \\ \\

    \theutterance \stepcounter{utterance}  
    & & \multicolumn{4}{p{0.6\linewidth}}{
        \cellcolor[rgb]{0.9,0.9,0.9}{
            \makecell[{{p{\linewidth}}}]{
                \texttt{\tiny{[GM$\rangle$P2]}}
                \texttt{Tá tú ag glacadh páirte i gcluiche tomhais comhoibríoch.} \\
\texttt{Is é an sprioc a fháil amach an bhfuil an greille seo agus pictiúr eile nach féidir liomsa amháin a fheiceáil, mar an gcéanna.} \\
\texttt{Déan cur síos ar do ghreille ar dtús, le do thoil. Ansin, tabharfaidh mé mo chur síos agus is féidir linn ceisteanna a chur ar a chéile faoi na greillí chun a dhéanamh amach an bhfuil siad mar an gcéanna. Anois cuir tús le do chur síos gairid ar an ngreille le "TUAIRISC:" agus an cur síos ina dhiaidh. Ná cuir aon rud eile leis. Seo é do ghreille:} \\
\texttt{▢ ▢ X ▢ ▢} \\
\texttt{▢ ▢ X ▢ ▢} \\
\texttt{▢ ▢ X ▢ ▢} \\
\texttt{▢ ▢ X ▢ ▢} \\
\texttt{X X X X X} \\
            }
        }
    }
    & \\ \\

    \theutterance \stepcounter{utterance}  
    & & & \multicolumn{4}{p{0.6\linewidth}}{
        \cellcolor[rgb]{0.9,0.9,0.9}{
            \makecell[{{p{\linewidth}}}]{
                \texttt{\tiny{[GM$\langle$P2]}}
                \texttt{TUAIRISC:} \\
\texttt{Greille 5x5. Lár{-}cholún (colún 3) agus lár{-}líne (líne 5) lán le X, cruthú cruachPhairc. Na cealla eile folamh.} \\
            }
        }
    }
     \\ \\

    \theutterance \stepcounter{utterance}  
    & & & \multicolumn{2}{p{0.3\linewidth}}{
        \cellcolor[rgb]{0.9,0.9,0.9}{
            \makecell[{{p{\linewidth}}}]{
                \texttt{\tiny{[GM$|$GM]}}
                \texttt{continue} \\
            }
        }
    }
    & & \\ \\

    \theutterance \stepcounter{utterance}  
    & & \multicolumn{4}{p{0.6\linewidth}}{
        \cellcolor[rgb]{0.9,0.9,0.9}{
            \makecell[{{p{\linewidth}}}]{
                \texttt{\tiny{[GM$\rangle$P2]}}
                \texttt{Seo é mo CHUR SÍOSTUAIRISC: Greille 5x5. Trídchló ghearrthóir (colún 3) agus líne bhunúsach (ráille 5) atá líonta le X, ag cruthú cruth T. Na cealla eile folamh.} \\
\texttt{Anois cuir ceist chun gnéithe nua de mo ghreille a fháil amach a d'fhéadfadh a bheith difriúil ó do ghreille. Tosaigh le "CEIST:" agus ná cuir aon rud eile leis.} \\
            }
        }
    }
    & \\ \\

    \theutterance \stepcounter{utterance}  
    & & & \multicolumn{4}{p{0.6\linewidth}}{
        \cellcolor[rgb]{0.9,0.9,0.9}{
            \makecell[{{p{\linewidth}}}]{
                \texttt{\tiny{[GM$\langle$P2]}}
                \texttt{CEIST: An bhfuil do chruachPhairc (líne 5) i gcónaí ar fad (colúin 1{-}5), nó an bhfuil sé teoranta go dtí an trídchló ghearrthóir (colún 3) amháin?} \\
            }
        }
    }
     \\ \\

    \theutterance \stepcounter{utterance}  
    & & & \multicolumn{2}{p{0.3\linewidth}}{
        \cellcolor[rgb]{0.9,0.9,0.9}{
            \makecell[{{p{\linewidth}}}]{
                \texttt{\tiny{[GM$|$GM]}}
                \texttt{continue} \\
            }
        }
    }
    & & \\ \\

    \theutterance \stepcounter{utterance}  
    & & \multicolumn{4}{p{0.6\linewidth}}{
        \cellcolor[rgb]{0.9,0.9,0.9}{
            \makecell[{{p{\linewidth}}}]{
                \texttt{\tiny{[P1$\langle$GM]}}
                \texttt{Seo é mo CHUR SÍOSTUAIRISC:} \\
\texttt{Greille 5x5. Lár{-}cholún (colún 3) agus lár{-}líne (líne 5) lán le X, cruthú cruachPhairc. Na cealla eile folamh.} \\
\texttt{CEIST: An bhfuil do chruachPhairc (líne 5) i gcónaí ar fad (colúin 1{-}5), nó an bhfuil sé teoranta go dtí an trídchló ghearrthóir (colún 3) amháin?Tosaigh do fhreagra le "FREAGRA:" agus ná cuir aon rud eile leis.} \\
            }
        }
    }
    & \\ \\

    \theutterance \stepcounter{utterance}  
    & \multicolumn{4}{p{0.6\linewidth}}{
        \cellcolor[rgb]{0.9,0.9,0.9}{
            \makecell[{{p{\linewidth}}}]{
                \texttt{\tiny{[P1$\rangle$GM]}}
                \texttt{FREAGRA: Tá mo líne bhunúsach (ráille 5) lán go hiomlán (colúin 1{-}5).} \\
            }
        }
    }
    & & \\ \\

    \theutterance \stepcounter{utterance}  
    & & & \multicolumn{2}{p{0.3\linewidth}}{
        \cellcolor[rgb]{0.9,0.9,0.9}{
            \makecell[{{p{\linewidth}}}]{
                \texttt{\tiny{[GM$|$GM]}}
                \texttt{continue} \\
            }
        }
    }
    & & \\ \\

    \theutterance \stepcounter{utterance}  
    & & \multicolumn{4}{p{0.6\linewidth}}{
        \cellcolor[rgb]{0.9,0.9,0.9}{
            \makecell[{{p{\linewidth}}}]{
                \texttt{\tiny{[P1$\langle$GM]}}
                \texttt{Anois cuir ceist chun gnéithe nua de mo ghreille a fháil amach a d'fhéadfadh a bheith difriúil ó do ghreille. Tosaigh le "CEIST:" agus ná cuir aon rud eile leis.} \\
            }
        }
    }
    & \\ \\

    \theutterance \stepcounter{utterance}  
    & \multicolumn{4}{p{0.6\linewidth}}{
        \cellcolor[rgb]{0.9,0.9,0.9}{
            \makecell[{{p{\linewidth}}}]{
                \texttt{\tiny{[P1$\rangle$GM]}}
                \texttt{CEIST: An bhfuil aon Xanna breise agat lasmuigh den trídchló ghearrthóir (colún 3) agus den líne bhunúsach (ráille 5), mar shampla i gcolún 1, 2, 4, nó 5 ar rfoill 1{-}4?} \\
            }
        }
    }
    & & \\ \\

    \theutterance \stepcounter{utterance}  
    & & & \multicolumn{2}{p{0.3\linewidth}}{
        \cellcolor[rgb]{0.9,0.9,0.9}{
            \makecell[{{p{\linewidth}}}]{
                \texttt{\tiny{[GM$|$GM]}}
                \texttt{continue} \\
            }
        }
    }
    & & \\ \\

    \theutterance \stepcounter{utterance}  
    & & \multicolumn{4}{p{0.6\linewidth}}{
        \cellcolor[rgb]{0.9,0.9,0.9}{
            \makecell[{{p{\linewidth}}}]{
                \texttt{\tiny{[GM$\rangle$P2]}}
                \texttt{FREAGRA: Tá mo líne bhunúsach (ráille 5) lán go hiomlán (colúin 1{-}5).} \\
\texttt{CEIST: An bhfuil aon Xanna breise agat lasmuigh den trídchló ghearrthóir (colún 3) agus den líne bhunúsach (ráille 5), mar shampla i gcolún 1, 2, 4, nó 5 ar rfoill 1{-}4?Tosaigh do fhreagra le "FREAGRA:" agus ná cuir aon rud eile leis.} \\
            }
        }
    }
    & \\ \\

    \theutterance \stepcounter{utterance}  
    & & & \multicolumn{4}{p{0.6\linewidth}}{
        \cellcolor[rgb]{0.9,0.9,0.9}{
            \makecell[{{p{\linewidth}}}]{
                \texttt{\tiny{[GM$\langle$P2]}}
                \texttt{FREAGRA: Níl, nílimid ach ag an trídchló ghearrthóir (colún 3, rfoill 1{-}4) agus an líne bhunúsach (ráille 5, colúin 1{-}5).} \\
\texttt{CEIST: An bhfuil an X i gcolún 3, raille 3 (an lár) agat?} \\
            }
        }
    }
     \\ \\

    \theutterance \stepcounter{utterance}  
    & & & \multicolumn{2}{p{0.3\linewidth}}{
        \cellcolor[rgb]{0.9,0.9,0.9}{
            \makecell[{{p{\linewidth}}}]{
                \texttt{\tiny{[GM$|$GM]}}
                \texttt{continue} \\
            }
        }
    }
    & & \\ \\

    \theutterance \stepcounter{utterance}  
    & & \multicolumn{4}{p{0.6\linewidth}}{
        \cellcolor[rgb]{0.9,0.9,0.9}{
            \makecell[{{p{\linewidth}}}]{
                \texttt{\tiny{[GM$\rangle$P2]}}
                \texttt{Anois cuir ceist chun gnéithe nua de mo ghreille a fháil amach a d'fhéadfadh a bheith difriúil ó do ghreille. Tosaigh le "CEIST:" agus ná cuir aon rud eile leis.} \\
            }
        }
    }
    & \\ \\

    \theutterance \stepcounter{utterance}  
    & & & \multicolumn{4}{p{0.6\linewidth}}{
        \cellcolor[rgb]{0.9,0.9,0.9}{
            \makecell[{{p{\linewidth}}}]{
                \texttt{\tiny{[GM$\langle$P2]}}
                \texttt{CEIST: An bhfuil an cruth agat díreach leis an dtrídchló ghearrthóir (colún 3, rfoill 1{-}5) agus an líne bhunúsach (ráille 5, colúin 1{-}5), gan aon Xanna eile?} \\
            }
        }
    }
     \\ \\

    \theutterance \stepcounter{utterance}  
    & & & \multicolumn{2}{p{0.3\linewidth}}{
        \cellcolor[rgb]{0.9,0.9,0.9}{
            \makecell[{{p{\linewidth}}}]{
                \texttt{\tiny{[GM$|$GM]}}
                \texttt{continue} \\
            }
        }
    }
    & & \\ \\

    \theutterance \stepcounter{utterance}  
    & & \multicolumn{4}{p{0.6\linewidth}}{
        \cellcolor[rgb]{0.9,0.9,0.9}{
            \makecell[{{p{\linewidth}}}]{
                \texttt{\tiny{[P1$\langle$GM]}}
                \texttt{FREAGRA: Níl, nílimid ach ag an trídchló ghearrthóir (colún 3, rfoill 1{-}4) agus an líne bhunúsach (ráille 5, colúin 1{-}5).} \\
\texttt{CEIST: An bhfuil an X i gcolún 3, raille 3 (an lár) agat?} \\
\texttt{CEIST: An bhfuil an cruth agat díreach leis an dtrídchló ghearrthóir (colún 3, rfoill 1{-}5) agus an líne bhunúsach (ráille 5, colúin 1{-}5), gan aon Xanna eile?Tosaigh do fhreagra le "FREAGRA:" agus ná cuir aon rud eile leis.} \\
            }
        }
    }
    & \\ \\

    \theutterance \stepcounter{utterance}  
    & \multicolumn{4}{p{0.6\linewidth}}{
        \cellcolor[rgb]{0.9,0.9,0.9}{
            \makecell[{{p{\linewidth}}}]{
                \texttt{\tiny{[P1$\rangle$GM]}}
                \texttt{FREAGRA: Tá, tá an X i gcolún 3, raille 3 agam. Tá, tá an cruth agam díreach leis an trídchló ghearrthóir (colún 3, rfoill 1{-}5) agus an líne bhunúsach (ráille 5, colúin 1{-}5), gan aon Xanna eile.} \\
            }
        }
    }
    & & \\ \\

    \theutterance \stepcounter{utterance}  
    & & & \multicolumn{2}{p{0.3\linewidth}}{
        \cellcolor[rgb]{0.9,0.9,0.9}{
            \makecell[{{p{\linewidth}}}]{
                \texttt{\tiny{[GM$|$GM]}}
                \texttt{continue} \\
            }
        }
    }
    & & \\ \\

    \theutterance \stepcounter{utterance}  
    & & \multicolumn{4}{p{0.6\linewidth}}{
        \cellcolor[rgb]{0.9,0.9,0.9}{
            \makecell[{{p{\linewidth}}}]{
                \texttt{\tiny{[P1$\langle$GM]}}
                \texttt{Anois cuir ceist chun gnéithe nua de mo ghreille a fháil amach a d'fhéadfadh a bheith difriúil ó do ghreille. Tosaigh le "CEIST:" agus ná cuir aon rud eile leis.} \\
            }
        }
    }
    & \\ \\

    \theutterance \stepcounter{utterance}  
    & \multicolumn{4}{p{0.6\linewidth}}{
        \cellcolor[rgb]{0.9,0.9,0.9}{
            \makecell[{{p{\linewidth}}}]{
                \texttt{\tiny{[P1$\rangle$GM]}}
                \texttt{CEIST: An bhfuil d'Xanna ina gcomharthaí "X" le fleisc (capital letters), nó an bhfuil siad scríofa mar "x" mall?} \\
            }
        }
    }
    & & \\ \\

    \theutterance \stepcounter{utterance}  
    & & & \multicolumn{2}{p{0.3\linewidth}}{
        \cellcolor[rgb]{0.9,0.9,0.9}{
            \makecell[{{p{\linewidth}}}]{
                \texttt{\tiny{[GM$|$GM]}}
                \texttt{continue} \\
            }
        }
    }
    & & \\ \\

    \theutterance \stepcounter{utterance}  
    & & \multicolumn{4}{p{0.6\linewidth}}{
        \cellcolor[rgb]{0.9,0.9,0.9}{
            \makecell[{{p{\linewidth}}}]{
                \texttt{\tiny{[GM$\rangle$P2]}}
                \texttt{FREAGRA: Tá, tá an X i gcolún 3, raille 3 agam. Tá, tá an cruth agam díreach leis an trídchló ghearrthóir (colún 3, rfoill 1{-}5) agus an líne bhunúsach (ráille 5, colúin 1{-}5), gan aon Xanna eile.} \\
\texttt{CEIST: An bhfuil d'Xanna ina gcomharthaí "X" le fleisc (capital letters), nó an bhfuil siad scríofa mar "x" mall?Tosaigh do fhreagra le "FREAGRA:" agus ná cuir aon rud eile leis.} \\
            }
        }
    }
    & \\ \\

    \theutterance \stepcounter{utterance}  
    & & & \multicolumn{4}{p{0.6\linewidth}}{
        \cellcolor[rgb]{0.9,0.9,0.9}{
            \makecell[{{p{\linewidth}}}]{
                \texttt{\tiny{[GM$\langle$P2]}}
                \texttt{FREAGRA: Is comharthaí "X" le fleisc (capital letters) iad.} \\
\texttt{CEIST: An 5x5 an méid a bhíonn sa ghreille agat, nó an bhfuil sí níos mó/níos lú?} \\
            }
        }
    }
     \\ \\

    \theutterance \stepcounter{utterance}  
    & & & \multicolumn{2}{p{0.3\linewidth}}{
        \cellcolor[rgb]{0.9,0.9,0.9}{
            \makecell[{{p{\linewidth}}}]{
                \texttt{\tiny{[GM$|$GM]}}
                \texttt{continue} \\
            }
        }
    }
    & & \\ \\

    \theutterance \stepcounter{utterance}  
    & & \multicolumn{4}{p{0.6\linewidth}}{
        \cellcolor[rgb]{0.9,0.9,0.9}{
            \makecell[{{p{\linewidth}}}]{
                \texttt{\tiny{[GM$\rangle$P2]}}
                \texttt{Anois cuir ceist chun gnéithe nua de mo ghreille a fháil amach a d'fhéadfadh a bheith difriúil ó do ghreille. Tosaigh le "CEIST:" agus ná cuir aon rud eile leis.} \\
            }
        }
    }
    & \\ \\

    \theutterance \stepcounter{utterance}  
    & & & \multicolumn{4}{p{0.6\linewidth}}{
        \cellcolor[rgb]{0.9,0.9,0.9}{
            \makecell[{{p{\linewidth}}}]{
                \texttt{\tiny{[GM$\langle$P2]}}
                \texttt{CEIST: An bhfuil na cealla folamh agat marthana ach spásanna béile, nó an bhfuil siad marcathe le comhartha eile (mar "▢" nó ".")?} \\
            }
        }
    }
     \\ \\

    \theutterance \stepcounter{utterance}  
    & & & \multicolumn{2}{p{0.3\linewidth}}{
        \cellcolor[rgb]{0.9,0.9,0.9}{
            \makecell[{{p{\linewidth}}}]{
                \texttt{\tiny{[GM$|$GM]}}
                \texttt{continue} \\
            }
        }
    }
    & & \\ \\

    \theutterance \stepcounter{utterance}  
    & & \multicolumn{4}{p{0.6\linewidth}}{
        \cellcolor[rgb]{0.9,0.9,0.9}{
            \makecell[{{p{\linewidth}}}]{
                \texttt{\tiny{[P1$\langle$GM]}}
                \texttt{FREAGRA: Is comharthaí "X" le fleisc (capital letters) iad.} \\
\texttt{CEIST: An 5x5 an méid a bhíonn sa ghreille agat, nó an bhfuil sí níos mó/níos lú?} \\
\texttt{CEIST: An bhfuil na cealla folamh agat marthana ach spásanna béile, nó an bhfuil siad marcathe le comhartha eile (mar "▢" nó ".")?Tosaigh do fhreagra le "FREAGRA:" agus ná cuir aon rud eile leis.} \\
            }
        }
    }
    & \\ \\

    \theutterance \stepcounter{utterance}  
    & \multicolumn{4}{p{0.6\linewidth}}{
        \cellcolor[rgb]{0.9,0.9,0.9}{
            \makecell[{{p{\linewidth}}}]{
                \texttt{\tiny{[P1$\rangle$GM]}}
                \texttt{FREAGRA: Is greille 5x5 í. Tá na cealla folamh agam mar "▢".} \\
            }
        }
    }
    & & \\ \\

    \theutterance \stepcounter{utterance}  
    & & & \multicolumn{2}{p{0.3\linewidth}}{
        \cellcolor[rgb]{0.9,0.9,0.9}{
            \makecell[{{p{\linewidth}}}]{
                \texttt{\tiny{[GM$|$GM]}}
                \texttt{continue} \\
            }
        }
    }
    & & \\ \\

    \theutterance \stepcounter{utterance}  
    & & \multicolumn{4}{p{0.6\linewidth}}{
        \cellcolor[rgb]{0.9,0.9,0.9}{
            \makecell[{{p{\linewidth}}}]{
                \texttt{\tiny{[P1$\langle$GM]}}
                \texttt{Anois cuir ceist chun gnéithe nua de mo ghreille a fháil amach a d'fhéadfadh a bheith difriúil ó do ghreille. Tosaigh le "CEIST:" agus ná cuir aon rud eile leis.} \\
            }
        }
    }
    & \\ \\

    \theutterance \stepcounter{utterance}  
    & \multicolumn{4}{p{0.6\linewidth}}{
        \cellcolor[rgb]{0.9,0.9,0.9}{
            \makecell[{{p{\linewidth}}}]{
                \texttt{\tiny{[P1$\rangle$GM]}}
                \texttt{CEIST: An bhfuil an trídchló ghearrthóir agat i gcolún 3 (an lár) nó i gcolún eile?} \\
            }
        }
    }
    & & \\ \\

    \theutterance \stepcounter{utterance}  
    & & & \multicolumn{2}{p{0.3\linewidth}}{
        \cellcolor[rgb]{0.9,0.9,0.9}{
            \makecell[{{p{\linewidth}}}]{
                \texttt{\tiny{[GM$|$GM]}}
                \texttt{continue} \\
            }
        }
    }
    & & \\ \\

    \theutterance \stepcounter{utterance}  
    & & \multicolumn{4}{p{0.6\linewidth}}{
        \cellcolor[rgb]{0.9,0.9,0.9}{
            \makecell[{{p{\linewidth}}}]{
                \texttt{\tiny{[GM$\rangle$P2]}}
                \texttt{FREAGRA: Is greille 5x5 í. Tá na cealla folamh agam mar "▢".} \\
\texttt{CEIST: An bhfuil an trídchló ghearrthóir agat i gcolún 3 (an lár) nó i gcolún eile?Tosaigh do fhreagra le "FREAGRA:" agus ná cuir aon rud eile leis.} \\
            }
        }
    }
    & \\ \\

    \theutterance \stepcounter{utterance}  
    & & & \multicolumn{4}{p{0.6\linewidth}}{
        \cellcolor[rgb]{0.9,0.9,0.9}{
            \makecell[{{p{\linewidth}}}]{
                \texttt{\tiny{[GM$\langle$P2]}}
                \texttt{FREAGRA: Tá an trídchló ghearrthóir i gcolún 3 (an lár).} \\
\texttt{CEIST: An bhfuil na "▢" i do gcur síos chun an ghreille a chruthú go pictiúril, nó an bhfuil siad mar thaganna comhaid (markup) atá feicthe agat?} \\
            }
        }
    }
     \\ \\

    \theutterance \stepcounter{utterance}  
    & & & \multicolumn{2}{p{0.3\linewidth}}{
        \cellcolor[rgb]{0.9,0.9,0.9}{
            \makecell[{{p{\linewidth}}}]{
                \texttt{\tiny{[GM$|$GM]}}
                \texttt{continue} \\
            }
        }
    }
    & & \\ \\

    \theutterance \stepcounter{utterance}  
    & & \multicolumn{4}{p{0.6\linewidth}}{
        \cellcolor[rgb]{0.9,0.9,0.9}{
            \makecell[{{p{\linewidth}}}]{
                \texttt{\tiny{[GM$\rangle$P2]}}
                \texttt{Anois déan cinneadh. Cad a cheapann tú: an bhfuil do ghreille agus an ghreille eile a thuairiscítear mar an gcéanna? Scríobh "CINNEADH: greille céanna" má cheapann tú gurb í an ghreille chéanna í nó "CINNEADH: greillí éagsúla" má cheapann tú gur greillí éagsúla iad. Ná cuir aon rud eile leis.} \\
            }
        }
    }
    & \\ \\

    \theutterance \stepcounter{utterance}  
    & & & \multicolumn{4}{p{0.6\linewidth}}{
        \cellcolor[rgb]{0.9,0.9,0.9}{
            \makecell[{{p{\linewidth}}}]{
                \texttt{\tiny{[GM$\langle$P2]}}
                \texttt{CINNEADH: greille céanna} \\
            }
        }
    }
     \\ \\

    \theutterance \stepcounter{utterance}  
    & & & \multicolumn{2}{p{0.3\linewidth}}{
        \cellcolor[rgb]{0.9,0.9,0.9}{
            \makecell[{{p{\linewidth}}}]{
                \texttt{\tiny{[GM$|$GM]}}
                \texttt{continue} \\
            }
        }
    }
    & & \\ \\

    \theutterance \stepcounter{utterance}  
    & & & \multicolumn{2}{p{0.3\linewidth}}{
        \cellcolor[rgb]{0.9,0.9,0.9}{
            \makecell[{{p{\linewidth}}}]{
                \texttt{\tiny{[GM$|$GM]}}
                \texttt{success} \\
            }
        }
    }
    & & \\ \\

    \theutterance \stepcounter{utterance}  
    & & \multicolumn{4}{p{0.6\linewidth}}{
        \cellcolor[rgb]{0.9,0.9,0.9}{
            \makecell[{{p{\linewidth}}}]{
                \texttt{\tiny{[P1$\langle$GM]}}
                \texttt{FREAGRA: Tá an trídchló ghearrthóir i gcolún 3 (an lár).} \\
\texttt{CEIST: An bhfuil na "▢" i do gcur síos chun an ghreille a chruthú go pictiúril, nó an bhfuil siad mar thaganna comhaid (markup) atá feicthe agat?} \\
\texttt{Anois déan cinneadh. Cad a cheapann tú: an bhfuil do ghreille agus an ghreille eile a thuairiscítear mar an gcéanna? Scríobh "CINNEADH: greille céanna" má cheapann tú gurb í an ghreille chéanna í nó "CINNEADH: greillí éagsúla" má cheapann tú gur greillí éagsúla iad. Ná cuir aon rud eile leis.} \\
            }
        }
    }
    & \\ \\

    \theutterance \stepcounter{utterance}  
    & \multicolumn{4}{p{0.6\linewidth}}{
        \cellcolor[rgb]{0.9,0.9,0.9}{
            \makecell[{{p{\linewidth}}}]{
                \texttt{\tiny{[P1$\rangle$GM]}}
                \texttt{CINNEADH: greille céanna} \\
            }
        }
    }
    & & \\ \\

    \theutterance \stepcounter{utterance}  
    & & & \multicolumn{2}{p{0.3\linewidth}}{
        \cellcolor[rgb]{0.9,0.9,0.9}{
            \makecell[{{p{\linewidth}}}]{
                \texttt{\tiny{[GM$|$GM]}}
                \texttt{continue} \\
            }
        }
    }
    & & \\ \\

    \theutterance \stepcounter{utterance}  
    & & & \multicolumn{2}{p{0.3\linewidth}}{
        \cellcolor[rgb]{0.9,0.9,0.9}{
            \makecell[{{p{\linewidth}}}]{
                \texttt{\tiny{[GM$|$GM]}}
                \texttt{success} \\
            }
        }
    }
    & & \\ \\

\end{supertabular}
}

\end{document}
